\documentclass[12pt,a4paper]{article}
\usepackage[utf8]{inputenc}
\usepackage[T1]{fontenc}
\usepackage[spanish]{babel}
\usepackage{geometry}
\usepackage{xcolor}
\usepackage{enumitem}
\usepackage{titlesec}

\geometry{left=2.54cm,right=2.54cm,top=2.54cm,bottom=2.54cm}
\definecolor{rojoUMB}{HTML}{8B0000}
\titleformat{\section}{\color{rojoUMB}\normalfont\large\bfseries}{\thesection}{0.7em}{}
\setlist[itemize]{leftmargin=1.2cm}
\setlist[enumerate]{leftmargin=1.2cm}

\title{\textbf{Ejemplos de Sesión Clínica}\\[0.2cm]\large Guía de Intervención Virtual de Apoyo a Fisioterapeutas}
\author{Angie Alejandra Amado Téllez}
\date{Abril de 2026}

\begin{document}

\maketitle

\noindent Este documento presenta ejemplos orientativos de uso clínico de la Guía de Intervención digital para pacientes con enfermedad de Parkinson en fase temprana. Los ejemplos no reemplazan el juicio profesional del fisioterapeuta y deben adaptarse según la valoración individual.

\section{Ejemplo 1. Sesión centrada en marcha y cueing rítmico}

\textbf{Objetivo de sesión:} mejorar regularidad del paso y reducir episodios de congelamiento de la marcha.

\begin{itemize}
    \item \textbf{Perfil clínico sugerido:} paciente en estadio II, con episodios leves de \textit{freezing of gait} y disminución de velocidad de marcha.
    \item \textbf{Valoración inicial:}
    \begin{itemize}
        \item Aplicar TUG Cognitivo.
        \item Registrar observaciones de entorno y factores de riesgo domiciliario.
        \item Complementar con FGA si el paciente tolera la prueba.
    \end{itemize}
    \item \textbf{Intervención sugerida:}
    \begin{enumerate}
        \item Calentamiento con marcha en línea recta y cambios de dirección lentos.
        \item Activación del metrónomo entre 55 y 70 BPM según tolerancia.
        \item Ejercicios de marcha con señal auditiva y visual, buscando sincronía del paso.
        \item Práctica de inicio y reinicio de la marcha ante pausas breves.
    \end{enumerate}
    \item \textbf{Seguimiento clínico:} registrar respuesta del paciente, estabilidad, fatiga y necesidad de apoyo externo.
\end{itemize}

\section{Ejemplo 2. Sesión centrada en equilibrio funcional}

\textbf{Objetivo de sesión:} fortalecer control postural y seguridad en transferencias.

\begin{itemize}
    \item \textbf{Perfil clínico sugerido:} paciente en estadio I o II, con antecedentes de inestabilidad al ponerse de pie o girar.
    \item \textbf{Valoración inicial:}
    \begin{itemize}
        \item Aplicar Escala de Berg (BBS).
        \item Registrar observaciones clínicas relevantes en equilibrio estático y dinámico.
    \end{itemize}
    \item \textbf{Intervención sugerida:}
    \begin{enumerate}
        \item Trabajo de sedente a bípedo con supervisión.
        \item Transferencias controladas con apoyo verbal.
        \item Alcances funcionales en distintas direcciones.
        \item Tareas de equilibrio con base de sustentación reducida, según tolerancia.
    \end{enumerate}
    \item \textbf{Seguimiento clínico:} registrar signos de inseguridad, necesidad de asistencia y evolución del desempeño funcional.
\end{itemize}

\section{Ejemplo 3. Sesión centrada en doble tarea}

\textbf{Objetivo de sesión:} favorecer integración cognitivo-motora durante la marcha.

\begin{itemize}
    \item \textbf{Perfil clínico sugerido:} paciente con disminución del rendimiento al caminar mientras habla, cuenta o responde indicaciones.
    \item \textbf{Valoración inicial:}
    \begin{itemize}
        \item Aplicar TUG Cognitivo.
        \item Registrar tiempo, errores y observaciones de atención dividida.
    \end{itemize}
    \item \textbf{Intervención sugerida:}
    \begin{enumerate}
        \item Marcha con conteo progresivo simple.
        \item Marcha con denominación de categorías sencillas.
        \item Marcha con cambios de ritmo apoyados por metrónomo.
        \item Pausas terapéuticas para control respiratorio y reorganización atencional.
    \end{enumerate}
    \item \textbf{Seguimiento clínico:} comparar estabilidad, ritmo y precisión cognitiva entre tareas simples y dobles.
\end{itemize}

\section{Recomendaciones generales}

\begin{itemize}
    \item Iniciar con tareas de baja complejidad y progresar de manera gradual.
    \item Ajustar el BPM del metrónomo según respuesta motora y seguridad del paciente.
    \item No utilizar la Guía de Intervención como sustituto de la valoración profesional.
    \item Guardar la valoración al final de la sesión para dejar trazabilidad clínica completa.
\end{itemize}

\end{document}